\documentclass[conference]{IEEEtran}

\usepackage[T1]{fontenc}
\usepackage[utf8]{inputenc}
\usepackage[brazilian]{babel}
\usepackage{booktabs}
\usepackage{hyperref}
\usepackage{amsmath}
\usepackage{url}

\title{BookPulse: Análise Comportamental e Exploração de Sentimento em Resenhas de Livros}

\author{
\IEEEauthorblockN{João Pedro Rodrigues Silva}
\IEEEauthorblockA{\textit{Campus V, Centro Federal de Educação Tecnológica de Minas Gerais}\\
Contagem, MG, Brasil\\
\href{mailto:jprs1308@gmail.com}{jprs1308@gmail.com}}
\and
\IEEEauthorblockN{Henrique de Freitas Araujo}
\IEEEauthorblockA{\textit{Campus V, Centro Federal de Educação Tecnológica de Minas Gerais}\\
Contagem, MG, Brasil\\
\href{mailto:ricogames@gmail.com}{ricogames@gmail.com}}
}

\begin{document}
\maketitle

\begin{abstract}
Este trabalho apresenta o BookPulse, um pipeline reprodutível para analisar
relações entre notas, texto, gêneros, popularidade e atividade de leitores em
resenhas do Goodreads. Uma amostra de 10.000 resenhas foi enriquecida com
metadados da Open Library e do Google Books, validada por critérios de
identificadores, títulos, autores e anos, e reduzida a 8.033 resenhas
associadas a livros com correspondência externa confiável. A análise textual
combina variáveis de tamanho, um baseline lexical interpretável e VADER como
inferência exploratória. Os resultados mostram associação positiva moderada
entre a nota e o score lexical normalizado (Pearson $r=0{,}182$; Spearman
$\rho=0{,}247$), mas não permitem afirmar causalidade nem validar o sentimento
sem anotações humanas. Como produto aplicado, o projeto inclui uma interface
que liga livros, palavras recorrentes e títulos com perfis semelhantes. O
artigo documenta também as limitações da amostra, dos métodos lexicais e da
recomendação baseada em similaridade agregada.
\end{abstract}

\begin{IEEEkeywords}
análise de sentimento, resenhas de livros, Goodreads, ciência de dados,
processamento de linguagem natural, sistemas de recomendação
\end{IEEEkeywords}

\section{Introdução}

Resenhas de livros combinam uma avaliação numérica com uma justificativa textual
que pode conter elogios, críticas, descrição de enredo, ironia e sentimentos
mistos. A nota atribuída pelo leitor não é equivalente à polaridade expressa
no texto. Esta distinção é central para evitar que uma análise classifique
automaticamente uma resenha como positiva apenas porque ela recebeu cinco
estrelas.

O objetivo deste trabalho é construir uma base confiável para estudar o
comportamento de leitores e produzir uma exploração visual dos padrões de
linguagem associados aos livros. As perguntas orientadoras são: (i) como
extensão, sentimento lexical, votos e comentários se associam à nota; (ii) que
diferenças aparecem entre gêneros e perfis de atividade; e (iii) como palavras
recorrentes podem ser utilizadas para encontrar livros com padrões textuais
semelhantes?

As contribuições são: um pipeline com validação de cruzamentos bibliográficos;
uma etapa textual auditável que preserva o texto original; uma especificação
de validação humana para futura avaliação supervisionada; e o BookPulse, uma
interface que visualiza palavras compartilhadas entre livros.

\section{Dados e escopo}

O dataset original contém aproximadamente 15 milhões de resenhas. Para tornar a
execução local e reproduzível, foi utilizada uma amostra de 10.000 registros
por reservoir sampling em uma passagem, com semente 42. A amostra não deve ser
interpretada como representação perfeita da população total.

Os livros foram enriquecidos no nível de $book\_id$. Open Library foi consultada
priorizando ISBN-13 e ISBN-10, seguida por buscas de título e autor. Google Books
foi usado como fonte independente de metadados. O validador combinou
concordância de ISBN, similaridade de título, autores, ano de publicação e
concordância entre fontes. Apenas os estados \texttt{MUITO\_ALTO},
\texttt{ALTO} e \texttt{MEDIO} foram utilizados na análise principal.

\begin{table}[!t]
\caption{Cobertura observada após a validação}
\label{tab:cobertura}
\centering
\begin{tabular}{lr}
\toprule
\textbf{Indicador} & \textbf{Quantidade}\\
\midrule
Reviews na amostra & 10.000\\
Livros únicos na amostra & 8.916\\
Livros com match confiável & 7.031\\
Reviews com match confiável & 8.033\\
Usuários na base final & 7.350\\
Reviews com duplicata textual exata & 66\\
\bottomrule
\end{tabular}
\end{table}

\section{Metodologia}

\subsection{Preparação e gêneros}

Duplicatas de reviews e livros foram removidas, números e datas foram
normalizados e o texto original foi preservado. Reviews vazias não aparecem na
base bruta final; 50 registros, contudo, não produzem tokens úteis no
refinamento lexical. Esse caso é diferente de ausência de texto e é reportado
separadamente.

Os gêneros foram inferidos combinando \texttt{popular\_shelves}, assuntos da
Open Library e categorias do Google Books. Evidências externas receberam maior
peso que shelves sociais, com no máximo três gêneros por livro. Foram atribuídos
ao menos um gênero a 6.971 dos 7.031 livros; 60 livros permaneceram como
\texttt{unclassified}.

\subsection{Representações textuais}

O texto original foi mantido para auditoria. Para a análise exploratória,
HTML, URLs e caracteres de controle foram tratados, e foram extraídos tokens,
quantidade de palavras, frases e razão de vocabulário único. O baseline lexical
usa um vocabulário pequeno e explícito em inglês. Para cada token, a polaridade
é $+1$, $-1$ ou zero; um negador nas três posições anteriores inverte a
polaridade. O score normalizado é:
\[
s_{100} = 100 \frac{N_{+}-N_{-}}{N_{\text{tokens}}}.
\]

Esse procedimento é interpretável, mas não modela adequadamente ironia,
composição semântica, contexto literário ou sentimentos mistos.

\subsection{VADER e validação proposta}

VADER foi incluído como segundo método exploratório, produzindo
\texttt{neg}, \texttt{neu}, \texttt{pos} e \texttt{compound}. Foram utilizados
os limiares convencionais $-0{,}05$ e $0{,}05$. O método foi desenvolvido para
linguagem informal de redes sociais; sua adequação ao domínio de livros não foi
presumida.

O projeto gera uma amostra estratificada, sem a nota original, para anotação
humana. As classes são \texttt{negative}, \texttt{neutral} e
\texttt{positive}; \texttt{mixed} e \texttt{unclear} podem ser registrados
durante a anotação. Com rótulos humanos, a etapa seguinte prevê a comparação
entre baseline, VADER e TF-IDF com unigramas/bigramas e regressão logística.
As métricas planejadas são macro-F1, F1 por classe, balanced accuracy, matriz
de confusão e Cohen's kappa entre anotadores. Como as anotações ainda não foram
preenchidas, nenhuma acurácia de classificador é reivindicada neste trabalho.

\subsection{Recomendação e visualização}

O site agrega as resenhas por livro. Para cada título, calcula proporções das
classes VADER, palavras recorrentes em todas as resenhas e palavras
recorrentes em subconjuntos positivo e negativo. A recomendação combina
gêneros compartilhados, proximidade do score de sentimento e quantidade de
reviews como critério secundário. O grafo apresenta o livro no centro, palavras
no primeiro anel e títulos que compartilham esses termos no anel externo.

Essa recomendação é uma prova de conceito exploratória, não uma avaliação
editorial nem um sistema validado por usuários. Capas são carregadas por URL
da Open Library a partir de ISBN, com fallback para tamanhos menores e imagem
gerada.

\section{Resultados}

\subsection{Perfil da base}

A base final contém 8.033 reviews, 7.031 livros e 7.350 usuários. Entre as
7.753 reviews com nota válida, a média foi 3,87, com mediana 4. Os dez
gêneros mais frequentes foram liderados por fantasy (1.032 reviews), romance
(935), science fiction (700), children (696) e mystery (659).

\subsection{Sentimento exploratório}

O VADER classificou 6.042 reviews (75,2\%) como positivas, 723 (9,0\%) como
neutras e 1.268 (15,8\%) como negativas. Esses valores são resultados de
inferência exploratória, não estimativas de prevalência validada.

\begin{table}[!t]
\caption{Nota e indicadores textuais}
\label{tab:nota}
\centering
\begin{tabular}{rrrr}
\toprule
\textbf{Nota} & \textbf{Reviews} & \textbf{Palavras} & \textbf{Compound}\\
\midrule
1 & 245 & 103,3 & -0,130\\
2 & 589 & 127,5 & 0,192\\
3 & 1.669 & 122,2 & 0,430\\
4 & 2.701 & 126,9 & 0,566\\
5 & 2.549 & 125,1 & 0,620\\
\bottomrule
\end{tabular}
\end{table}

O score lexical normalizado apresentou correlação de Pearson $r=0{,}182$ e
Spearman $\rho=0{,}247$ com a nota. Para comparação, o tamanho da review teve
Pearson $r=0{,}012$ e Spearman $\rho=-0{,}015$. Votos e comentários também
apresentaram associações pequenas. A diferença entre os coeficientes lexical e
VADER e a ausência de validação humana impedem tratar esses valores como prova
de que os métodos identificaram corretamente o sentimento.

\section{Discussão}

Os resultados são coerentes com a hipótese de que notas mais altas tendem a
coocorrer com linguagem mais positiva, mas a associação é moderada e não
elimina casos de crítica pontual, elogio com ressalvas ou ironia. O baixo
relacionamento entre tamanho e nota sugere que extensão, isoladamente, não é
um indicador suficiente da avaliação do leitor.

Há dependência entre observações: um usuário pode escrever várias reviews e um
livro pode receber vários textos. Além disso, 66 reviews possuem texto
exatamente duplicado. Uma etapa supervisionada futura deve evitar vazamento
com divisões agrupadas por usuário ou livro e deve manter um teste final
intocado. O trabalho também não permite conclusões causais sobre votos,
comentários, popularidade ou sentimento.

\section{Limitações e próximos passos}

O principal limite atual é a ausência de anotações humanas. O vocabulário
lexical é pequeno e VADER não foi criado especificamente para resenhas
literárias. A amostra é pequena diante da população original e os matches
externos, embora filtrados, dependem de metadados de edição. A recomendação
também não foi avaliada por métricas de ranking ou estudo com usuários.

Como próximos passos, recomenda-se: (i) concluir a anotação independente por
dois anotadores; (ii) calcular kappa e métricas por classe; (iii) comparar
VADER, TF-IDF e um transformer com splits por usuário e livro; (iv) avaliar
calibração e intervalos de confiança; e (v) realizar uma avaliação de
recomendações com participantes ou histórico de cliques.

\section{Conclusão}

O BookPulse entrega um pipeline coeso entre amostragem, enriquecimento,
validação, análise textual e visualização. A base final permite observar
associações entre nota e linguagem, além de explorar conexões entre livros por
palavras compartilhadas. O resultado mais importante é metodológico: o projeto
mantém a distinção entre nota, sentimento e predição, apresenta o baseline como
exploratório e explicita a validação humana ainda necessária para conclusões
de desempenho.

\begin{thebibliography}{00}
\bibitem{pang2002} B. Pang, L. Lee e S. Vaithyanathan, ``Thumbs up? Sentiment
Classification using Machine Learning Techniques,'' \textit{Proc. ACL}, 2002,
doi: 10.3115/1118693.1118704.
\bibitem{hutto2014} C. J. Hutto e E. Gilbert, ``VADER: A Parsimonious
Rule-based Model for Sentiment Analysis of Social Media Text,'' \textit{Proc.
ICWSM}, 2014, doi: 10.1609/icwsm.v8i1.14550.
\bibitem{socher2013} R. Socher et al., ``Recursive Deep Models for Semantic
Compositionality Over a Sentiment Treebank,'' \textit{Proc. EMNLP}, 2013.
\bibitem{devlin2019} J. Devlin et al., ``BERT: Pre-training of Deep
Bidirectional Transformers for Language Understanding,'' \textit{Proc. NAACL},
2019, doi: 10.18653/v1/N19-1423.
\bibitem{sklearn} F. Pedregosa et al., ``Scikit-learn: Machine Learning in
Python,'' \textit{JMLR}, vol. 12, pp. 2825--2830, 2011.
\bibitem{huggingface} Hugging Face, ``Text classification,'' 2024. [Online].
Available: \url{https://huggingface.co/docs/transformers/en/tasks/sequence_classification}
\end{thebibliography}

\end{document}
